\subsection{A Note on Irrational Exponents} 
We have now defined what it means to write any \makeblank as an exponent.
\begin{definition}
For positive integers $n$ and any real number $x$,
$$x^n=\underbrace{x\cdot x\cdots x}_{n\ \mathrm{copies}}$$
For any real number $x\neq 0$,
$$x^0=1$$
For any negative integer $-n$ and real number $x\neq 0$,
$$x^{-n}=\dfrac{1}{x^n}$$
For any rational number $\frac{p}{q}$ written in lowest terms and any real number $x$ for which the expressions on the right are defined,
$$x^{1/q}=\sqrt[q]{x}
\qquad
x^{p/q}=\sqrt[q]{x^p}=\left(\sqrt[q]{x}\right)^p$$
\end{definition}
We notice that $x^y$ is defined for all rational $y$ only when \makeblank[1in]. 